\section{Lessons, open points and conclusion}
\label{sec:lessons}

\subsection{Lessons}
\begin{enumerate}[leftmargin=1.4em]
  \item \textbf{Measure inside the kernel.} The AXI counters (read waits, bursts, latency) pointed at every
        improvement: the strided A reads, then the per-tile reloads. Without them the PE 32\x8 result would have
        looked like a bandwidth problem, which it never was (\cref{sec:eval:roofline}).
  \item \textbf{The HLS clock estimate is not a limit.} HLS deepens pipelines for a tighter target even when its
        estimate stays the same; only a routed build decides. An earlier conclusion that 225\,MHz was not worth a
        build was wrong for this reason.
  \item \textbf{Try the Vivado strategies before changing the kernel.} At 225\,MHz the difference between failing and
        passing was the physical optimisation of a few high-fanout nets.
  \item \textbf{Check the whole design's timing.} The most marginal path in every build is inside the PCIe hard block,
        not in the kernel; a flow that only looked at the kernel clock would have loaded a failing build.
  \item \textbf{Silent wrong answers are the worst failure.} Sizes that are not whole bus words, and the order of
        bring-up steps (calibrating DDR before the driver's probe resets the SoC), both produced wrong results without
        errors. Both are now checked or enforced by the tools.
  \item \textbf{Vendor tools need guards.} \code{vitis-run} exits 0 on a failed synthesis; Vivado removes logic from
        constant FIFOs and then fails routing (\code{Route 35-13}); old hlslib constructs fail in a new HLS. Each
        needed a workaround in the flow or the RTL.
  \item \textbf{Traceability pays off.} Because every row carries the build id, the configuration hash and the git
        state, two real bugs in the result handling were found and fixed afterwards: a run-row id that broke
        \code{--resume}, and kernel utilisation columns read one position off.
  \item \textbf{Check exactly, with data that switches.} A relative tolerance and small float values hid DDR3 read
        errors that exact integer checks over the full code range exposed (\cref{sec:formats:ddr}). Timing met in
        STA is not proof that a design computes correctly on the card.
  \item \textbf{Validate the kernel's own assumptions.} Two tile rules in \code{gemm\_hls} were implicit; a build
        that broke one met timing and produced wrong results, then hung. A ten-case g++ simulation confirmed both
        rules in seconds; they are now validation errors (\cref{sec:formats:tiles}).
  \item \textbf{Narrow formats move the limit.} Without the DSP bottleneck, the C-tile BRAM, the LUTs and the fan-out
        of the HLS dataflow control decide how large and how fast an array can be (\cref{sec:formats:wide}).
\end{enumerate}

\subsection{Open points}
\begin{itemize}
  \item \textbf{H2C throughput} (\SI{0.82}{GB/s}): the generated driver's 1\,ms DMA stop. Options: a shorter delay in
        the driver copy, or a bridge that back-pressures the DMA reader after \code{length} bytes so that nothing is in
        flight at the stop. Both touch the class of PCIe fault that froze the host early on, so they need a decision
        and careful testing.
  \item \textbf{PCIe core timing}: build the PHY as Gen2 x4 (the slot's width here), or keep x8 and rebuild with
        another preset when only that path fails.
  \item \textbf{Clock beyond 225\,MHz} (180\,MHz for 64 PEs) needs kernel changes (registering the dataflow control
        and the C-pipe enables).
  \item \textbf{A wider accumulator} (e.g.\ int8 inputs, int32 sums) and power-of-two weights (DenseShift) need
        separate A/B/C types in the kernel; the flow's configuration and result columns are prepared for them.
  \item \textbf{DDR3 read margin}: the host-side tuning works, but a permanent fix belongs in the PHY (per-bit
        delays, DQS capture) or in a calibration under load; either rebuilds every bitstream.
  \item \textbf{The fixed loss per k step} (\cref{sec:formats:kstep}) is not yet located in the kernel.
  \item \textbf{Memory map at PE 32\x8} was not measured; it should matter more there than at PE 8\x8.
  \item \textbf{Overlap of transfers and compute} (double-buffering matrices in DDR) would hide most of the transfer
        time for streams of GEMMs.
\end{itemize}

\subsection{Conclusion}
The project delivers a working, measured GEMM accelerator on a card without a vendor shell: a CPU-less LiteX SoC
with PCIe DMA into two DDR3 channels, an adapted \code{gemm\_hls} kernel with performance counters, a C++ host
runtime, and a JSON-driven flow whose DSE driver built, loaded and benchmarked twenty configurations unattended. The
best configuration, a 32\x8 PE array with transposed A, 1024\x512 tiles and a 225\,MHz clock closed by Vivado's
\code{aggressive} strategy, computes float32 GEMM at \textbf{\gflops{110.9}}, \SI{96.2}{\percent} of its peak, with
every result verified and traceable to the bitstream and source that produced it. With the element format as one
more configuration key, the same kernel source computes int8 GEMM at \SI{426.6}{GOP/s} and int4 at
\textbf{\SI{1370.9}{GOP/s}}, exactly checked on the card.
